
\section{一个玩具模型：有理拉格朗日子流形}
\label{sec:6-2}
\liftIndexRangeStart{I364}

如上一节所述，系数环
$\Lambda_{0,\mathrm{nov}}$ 不是诺特环。
因此，即使定义出谱序列，
也不易证明 $E_r^{p,q}$ 项有限生成，
或证明 (D.2) 收敛。
定理 D 和定理 6.1.4 的证明，
主要就在于解决这些问题。
在构造一般情形的谱序列之前，
我们先在一个简单情形中
展示其构造及其在定理 I 证明中的用法。
这个情形的系数环是诺特环。

我们以能量滤过构造谱序列。
先回顾一些记号。
设 $L$ 是 $(M,\omega)$ 中
相对 spin 的拉格朗日子流形，
$\psi$ 是哈密顿微分同胚。
假设 $L$ 无障碍，
并与 $\psi(L)$ 横截。
回忆定义 2.4.17 中的
$\Pi(L)=\pi_2(M,L)/\sim$。
它是覆叠
\[
\widetilde\Omega(L,\psi(L);\ell_0)
\longrightarrow
\Omega(L,\psi(L);\ell_0)
\]
的覆盖变换群
（见第 2.4.3 小节）。
取 $\ell_0=\ell_{\mathrm{can}}$，
右端已在第 2.3.2 小节定义。
对
$[(\ell,w)]\in
\widetilde\Omega(L,\psi(L);\ell_0)$，
令
$\mathcal A([(\ell,w)])=\int w^*\omega$。
于是有称为\emph{能量}的群同态
$E:\Pi(L)\to\mathbb R$，定义为
\[
E(g)=\mathcal A(g\cdot[(\ell,w)])
-\mathcal A([(\ell,w)]).
\]
它与 $[(\ell,w)]$ 的选取无关
（见第 2.2.1 小节）。
注意，第 2.2.1 小节是在
$L^{(0)}$ 与 $L^{(1)}$ 横截的情形
定义 $E:\Pi(L)\to\mathbb R$；
干净相交时的定义类似。

本节对 $L$ 加以下条件。

\begin{definition}{6.2.1}\label{def:6-2-1}
\begin{enumerate}[label=\textup{(6.2.2.\arabic*)},leftmargin=5.8em]
\item 若子群
\[
E(\Pi(L))
=\{E(\beta)\mid\beta\in\Pi(L)\}
\subset\mathbb R
\]
是离散的，
则称拉格朗日子流形 $L$ 为\emph{有理的}。
\item 设 $L\subset M$ 是有理拉格朗日子流形\liftIndex{I220}。
若存在满足条件 3.1.6 的子幺半群 $G$，
包含所有序对
$\{(E(\beta),\mu_L(\beta))
\mid\beta\in\Pi(L)\}$，
且满足以下两条，
则称 $L$ \emph{有理无障碍}\liftIndex{I221}：
\translatornote{原书实体页 372 的
(6.2.2.2) 要求 $G$ 包含全部
$\Pi(L)$ 的能量--马斯洛夫对，
但条件 3.1.6 要求 $G$ 的能量非负；
若 $E(\Pi(L))$ 含非零值，
群的取负运算使这两者按字面冲突。
同处及实体页 373 又把含单位的
$\pi(G)$ 写为 $\mathbb Z_{>0}$，
而该集合不含加法单位零。
这里保留原文，见疑误 A-0358。}
\begin{enumerate}[label=\textup{(6.2.3.\arabic*)},leftmargin=5.8em]
\item $\pi(G)\subset\mathbb R$
作为幺半群同构于 $\mathbb Z_{>0}$；
\item 按定义 3.6.16，
$(C(L),\mathfrak m)$ 在 $G$ 上无障碍。
\end{enumerate}
\end{enumerate}
\end{definition}

有理无障碍的弱版本、
以及含 bulk 形变的版本，
可类似定义。
对于后者，还须假设
$\boldsymbol b\in H^2(M;\Lambda_{0,\mathrm{nov}})$
可写成
\[
\boldsymbol b
=\sum_{(\lambda_i,\mu_i)\in G}
b_iT^{\lambda_i}e^{\mu_i/2},
\qquad b_i\in H^*(M;R).
\]

\begin{remark}{6.2.4}\label{rem:6-2-4}
另有文献，例如 \textbf{[Fuk03I]}，
以不同方式定义 $L$ 的有理性。
假设 $[\omega]\in H^2(M;\mathbb Z)$。
取复线丛 $\mathcal L\to M$，
配备 $U(1)$-联络 $\nabla$，
使其第一陈形式为 $\omega$。
若 $\nabla$ 限制到 $L$ 后的和乐群
为有限阶，便称 $L$ 有理。
当 $M$ 单连通时，
此定义与定义 6.2.1 等价；
一般情形下则比定义 6.2.1 更严格。
\end{remark}
\liftIndexRangeEnd{I364}

令 $G$ 是定义 6.2.1 所给的子幺半群；
本节始终固定 $G$。
引入如下诺维科夫环：
\begin{equation}
\widehat\Lambda(G)
=\left\{
\sum_i a_iT^{\lambda_i}e^{n_i}
\in\Lambda_{\mathrm{nov}}
\ \middle|\
\lambda_i\in E(G),\ a_i\in R
\right\},
\tag{6.2.5.1}\label{eq:6-2-5-1}
\end{equation}
及其子环
\begin{equation}
\begin{aligned}
\widehat\Lambda_+(G)
&=\left\{
\sum_i a_iT^{\lambda_i}e^{n_i}
\in\Lambda_{0,\mathrm{nov}}
\ \middle|\
\lambda_i\in E(G)
\right\}\\
&=\Lambda_{0,\mathrm{nov}}
\cap\widehat\Lambda(G).
\end{aligned}
\tag{6.2.5.2}\label{eq:6-2-5-2}
\end{equation}
它们与定义 2.4.17 之后的环
$\Lambda^R(L)$ 有关，但并不相同。

$\widehat\Lambda(G)$ 中元素的次数
与其在 $\Lambda_{\mathrm{nov}}$ 中的次数相同。
诺维科夫环 $\widehat\Lambda(G)$
是 $\widehat\Lambda_+(G)$ 的局部化。
若 $G$ 满足 (6.2.3.1)，
则 $\widehat\Lambda_+(G)$ 是诺特环。
事实上，若 $R$ 为域，
其零次部分
$\widehat\Lambda_+(G)^{(0)}$
是离散赋值环。
这使谱序列的研究较为简便。
本节在 $L$ 有理且有理无障碍时，
以 $\widehat\Lambda(G)$ 上的谱序列证明定理 I。
此时，谱序列的收敛性
由标准论证得到。
另须指出，
这里要用的性质 (D.3)
不依赖于系数环的诺特性
（见第 6.4.2 小节）。

由 (6.2.3.1)，有
$E(G)=\mathbb Z_{>0}\cdot\lambda_0$
（$\lambda_0>0$），从而
\begin{equation}
\widehat\Lambda_+(G)
\cong R[[T^{\lambda_0}]][e,e^{-1}].
\tag{6.2.6}\label{eq:6-2-6}
\end{equation}
由 (6.1.1) 定义的
$\Lambda_{0,\mathrm{nov}}$ 上的滤过
在 $\widehat\Lambda_+(G)$ 上诱导滤过；令
\[
\mathfrak F^m\widehat\Lambda_+(G)
=F^{m\lambda_0}\widehat\Lambda_+(G).
\]
注意，$\Lambda_{0,\mathrm{nov}}$
作为 $\widehat\Lambda_+(G)$-模是平坦的，
因而诱导出定理 D 中的谱序列。
为在当前情形证明定理 I，
如上一节所述，
我们改在 $\widehat\Lambda(G)$
而非 $\widehat\Lambda_+(G)$ 上工作。
这时有
\begin{equation}
\widehat\Lambda(G)
\cong R[[T^{\lambda_0}]]
[T^{-\lambda_0}][e,e^{-1}].
\tag{6.2.7}\label{eq:6-2-7}
\end{equation}
式 (6.1.2) 中
$\Lambda_{\mathrm{nov}}$ 的滤过
也在 $\widehat\Lambda(G)$ 上诱导滤过；令
\[
\mathfrak F^m\widehat\Lambda(G)
=F^{m\lambda_0}\widehat\Lambda(G).
\]
于是
\[
\operatorname{gr}_q
(\mathfrak F\widehat\Lambda(G))
\cong R[e,e^{-1}].
\]

在由 $L$ 上光滑奇异
$(n-k)$-单形生成的自由 $R$-模
$S^k(L;R)$ 中，
取一个可数生成子模 $C^k(L;R)$。
用 $C(L;\widehat\Lambda(G))$
表示
$C(L;R)\otimes_R\widehat\Lambda(G)$
的完备化。
$\widehat\Lambda(G)$ 的滤过
在它上面诱导出
$\mathfrak F^*C(L;\widehat\Lambda(G))$；
总次数是 $C(L;R)$ 与
$\widehat\Lambda(G)$ 的次数之和。
因为 $L$ 在 $G$ 上无障碍，
存在定界余链 $b\in\mathcal M(L)$，
并可选取 $b$ 使其属于
$C(L;\widehat\Lambda(G))$。
由命题 3.6.10，
得到形变 $A_\infty$ 代数
$(C(L;\widehat\Lambda(G)),\mathfrak m^b)$，
其中 $\mathfrak m_0^b=0$。
于是 $\mathfrak m_1^b$
定义 $C(L;\widehat\Lambda(G))$ 上的上边界算子，
并给出以此环为系数的形变弗洛尔上同调
$HF((L,b),(L,b);\widehat\Lambda(G))$。
本节在不致混淆时，
省略记号中的 $b$。
$C(L;\widehat\Lambda(G))$ 上的滤过
诱导我们要构造的谱序列。

首先有作为
$\operatorname{gr}_*(\mathfrak F\widehat\Lambda(G))$-模
的同构\translatornote{原书实体页 373
在刚定义 $\mathfrak F\widehat\Lambda(G)$ 后，
把其伴随分次写成
$\operatorname{gr}_q(\mathfrak F\widehat\Lambda_+(G))$；
式 (6.2.8) 左端又写作
$C(L;\widehat\Lambda(L))$。
这里依同页已定义的滤过与复形
统一用 $G$ 且去掉前一处多出的 $+$；
见疑误 A-0359。}
\begin{equation}
\operatorname{gr}_*
(\mathfrak F C(L;\widehat\Lambda(G)))
\cong C(L;R)\otimes
\operatorname{gr}_*
(\mathfrak F\widehat\Lambda(G)).
\tag{6.2.8}\label{eq:6-2-8}
\end{equation}
由微分的构造（定义 3.5.6），
(6.2.8) 上诱导的微分是
$(-1)^n\partial\otimes\mathrm{id}$，
其中 $\partial$ 是 $C(L;R)$
上的经典边界算子。
因此得到谱序列 $E_r^{p,q}$，
其 $E_2$ 项为
\begin{equation}
E_2^{p,q}
=\bigoplus_{k=0}^{n}
H^k(L;R)\otimes
\operatorname{gr}_q
\bigl(\mathfrak F
\widehat\Lambda(G)^{(p-k)}\bigr).
\tag{6.2.9}\label{eq:6-2-9}
\end{equation}
这里 $\widehat\Lambda(G)^{(p-k)}$
表示 $p-k$ 次部分。
还可证明，
$HF^p(L,L;\widehat\Lambda(G))$
上存在滤过，使
\begin{equation}
E_\infty^{p,q}
\cong
\frac{\mathfrak F^q
HF^p(L,L;\widehat\Lambda(G))}
{\mathfrak F^{q+1}
HF^p(L,L;\widehat\Lambda(G))}.
\tag{6.2.10}\label{eq:6-2-10}
\end{equation}
上述谱序列的收敛性证明
适用于任意交换环 $R$。

注意，分别乘以
$T^{\lambda_0}$ 和 $e$
给出如下同构：
\[
\begin{aligned}
T^{\lambda_0}\cdot:
&\ \mathfrak F^qHF^p(L,L;\widehat\Lambda(G))
\longrightarrow
\mathfrak F^{q+1}HF^p(L,L;\widehat\Lambda(G)),\\
e\cdot:
&\ \mathfrak F^qHF^p(L,L;\widehat\Lambda(G))
\longrightarrow
\mathfrak F^qHF^{p+2}(L,L;\widehat\Lambda(G)).
\end{aligned}
\]
现在考虑带滤过 $\mathfrak F^*D$
的分次 $\widehat\Lambda(G)$-模 $D$，
其滤过与 $\widehat\Lambda(G)$ 的滤过相容，
即
\[
\mathfrak F^\ell\widehat\Lambda(G)
\cdot\mathfrak F^kD
\subset\mathfrak F^{k+\ell}D.
\]
记 $\widehat\Lambda(G)^{(0)}$
为 $\widehat\Lambda(G)$ 的零次部分。
由
\[
\widehat\Lambda(G)^{(0)}
\cong R[[T^{\lambda_0}]][T^{-\lambda_0}]
\]
可知，若 $R$ 为域，
$\widehat\Lambda(G)^{(0)}$ 也是域，
故 $D^p$ 是其上的向量空间。

\begin{definition}{6.2.11}\label{def:6-2-11}
定义 $t$-秩\liftIndex{I222}
（拓扑秩，或暂定秩）为
\[
\operatorname{t-rank}D
=\dim_{\widehat\Lambda(G)^{(0)}}D^p
+\dim_{\widehat\Lambda(G)^{(0)}}D^{p+1}.
\]
\end{definition}

由于乘以 $e$ 给出
$D^p\cong D^{p+2}$，
右端与 $p$ 无关。
以下令 $R=\mathbb Q$。
若弗洛尔上同调可在
$\Lambda_{0,\mathrm{nov}}(R)$ 上定义，
且 $L$ 有理，
本节剩余论证也适用于任意域 $R$。

\begin{lemma}{6.2.12}\label{lem:6-2-12}
对任意 $q$，有
\[
\begin{aligned}
&\operatorname{t-rank}
HF(L,L;\widehat\Lambda(G))\\
&=\dim_{\mathbb Q}
\frac{\mathfrak F^qHF^p(L,L;\widehat\Lambda(G))}
{\mathfrak F^{q+1}HF^p(L,L;\widehat\Lambda(G))}\\
&\quad+\dim_{\mathbb Q}
\frac{\mathfrak F^qHF^{p+1}(L,L;\widehat\Lambda(G))}
{\mathfrak F^{q+1}HF^{p+1}(L,L;\widehat\Lambda(G))}\\
&=\dim_{\mathbb Q}E_\infty^{p,q}
+\dim_{\mathbb Q}E_\infty^{p+1,q}.
\end{aligned}
\]
\end{lemma}

证明容易，故略。
下一个引理由定理 3.7.21
与定理 4.1.3 直接推出。

\begin{lemma}{6.2.13}\label{lem:6-2-13}
若 $\psi$ 是哈密顿微分同胚，
且 $L$ 与 $\psi(L)$ 横截，
则对任意 $b\in\mathcal M(L)$，
\[
\#(L\cap\psi(L))
\ge
\operatorname{t-rank}
HF((L,b),(L,b);\widehat\Lambda(G)).
\]
\end{lemma}

现在从下方估计
$\dim_{\mathbb Q}E_\infty^{p,q}
+\dim_{\mathbb Q}E_\infty^{p+1,q}$。
令 $\delta_r:E_r\to E_r$
为谱序列的上边界算子。
有 $\delta_r\delta_r=0$，且
\[
E_{r+1}=\ker\delta_r/\operatorname{Im}\delta_r.
\]
$E_r$ 是
$\operatorname{gr}_*(\mathfrak F\widehat\Lambda(G))$-模，
而原书写
$\operatorname{gr}_*(\mathfrak F\widehat\Lambda(G))
\cong\widehat\Lambda(G)$；
因而原书据定义 6.2.11
给 $E_r$ 定义 $t$-秩。
\translatornote{原书实体页 374
在此断言伴随分次环
$\operatorname{gr}_*(\mathfrak F\widehat\Lambda(G))$
与完成环 $\widehat\Lambda(G)$ 同构。
但由 (6.2.7) 的形式，
前者仅含能量方向的有限 Laurent 和，
后者还含无穷正向幂级数；
按字面不能同构。
后续 $t$-秩的移用因而需要额外论证；
见疑误 A-0360。}

\begin{lemma}{6.2.14}\label{lem:6-2-14}
\[
\operatorname{t-rank}E_r
-2\operatorname{t-rank}
\operatorname{Im}\delta_r
=\operatorname{t-rank}E_{r+1}
\le\operatorname{t-rank}E_r.
\]
\end{lemma}

\begin{proof}
存在如下 $\widehat\Lambda(G)$-模短正合列：
\[
\begin{aligned}
0&\longrightarrow\ker\delta_r
\longrightarrow E_r
\longrightarrow(\operatorname{Im}\delta_r)[-1]
\longrightarrow0,\\
0&\longrightarrow\operatorname{Im}\delta_r
\longrightarrow\ker\delta_r
\longrightarrow E_{r+1}
\longrightarrow0.
\end{aligned}
\]
由第一列得到
\[
\operatorname{t-rank}E_r
=\operatorname{t-rank}\ker\delta_r
+\operatorname{t-rank}\operatorname{Im}\delta_r.
\]
第二列遂推出引理。
\end{proof}

下面证明定理 I，
设 $L$ 在 $G$ 上经 bulk 形变后
为（弱）无障碍。
由 (6.2.9) 及 $\deg e=2$，
得到\translatornote{原书实体页 375
首式右侧写
$\sum t\text{-rank}_{\mathbb Q}H^*(L;\mathbb Q)$，
但定义 6.2.11 未给普通有理向量空间
定义这一记号；同页下一式明确以
有理上同调各次数的普通秩求和。
此处按后式排印，见疑误 A-0361。}
\[
\operatorname{t-rank}E_2
=\sum_k\operatorname{rank}_{\mathbb Q}
H^k(L;\mathbb Q).
\]
按 (D.3) 的记号，令
\[
\begin{aligned}
K_2
&=\operatorname{PD}
\bigl(\ker(H_*(L;\mathbb Q)
\longrightarrow H_*(M;\mathbb Q))\bigr)
\otimes_{\mathbb Q}
\operatorname{gr}_*
(\mathfrak F\widehat\Lambda(G)),\\
K_{r+1}
&=\frac{K_r\cap\ker(\delta_r)}
{K_r\cap\operatorname{Im}(\delta_r)},
\end{aligned}
\]
将在第 6.4.2 小节证明：
利用算子 $\mathfrak p$ 可得
\begin{equation}
\operatorname{Im}(\delta_r)\subset K_r.
\tag{6.2.15}\label{eq:6-2-15}
\end{equation}
故
\[
\operatorname{t-rank}K_{r+1}
\le\operatorname{t-rank}K_r
-\operatorname{t-rank}
\operatorname{Im}(\delta_r).
\]
结合引理 6.2.14，得到
\[
\begin{aligned}
\operatorname{t-rank}E_r
&=\operatorname{t-rank}E_2
-2\sum_{2\le i<r}
\operatorname{t-rank}
\operatorname{Im}(\delta_i)\\
&=\sum_k\operatorname{rank}_{\mathbb Q}
H^k(L;\mathbb Q)
-2\sum_{2\le i<r}
\operatorname{t-rank}
\operatorname{Im}(\delta_i)\\
&\ge\sum_k\operatorname{rank}_{\mathbb Q}
H^k(L;\mathbb Q)
-2\bigl(\operatorname{t-rank}K_2
-\operatorname{t-rank}K_r\bigr).
\end{aligned}
\]
因此
\[
\begin{aligned}
\operatorname{t-rank}
HF(L,L;\widehat\Lambda(G))
&=\operatorname{t-rank}E_\infty\\
&\ge\sum_k\operatorname{rank}_{\mathbb Q}
H^k(L;\mathbb Q)
-2\operatorname{t-rank}K_2.
\end{aligned}
\]
对于有理无障碍的拉格朗日子流形，
定理 I 的证明至此完成，
惟仍需证明 (6.2.15)。
\hfill$\square$
